\section{流态管理：让不同的时钟一起工作}

工厂、医院和交通网络都可能在局部指标改善时，让接口变得更难管理。每个环节都更快，并不保证整条链更稳。接下来的问题是：怎样让响应速度不同的过程持续协作？

\subsection{一个系统里的多只时钟}

符号可以迅速复制，变压器的交付需要更长时间，某些临床后果则要长期观察。速度差属于具体过程。神经信号很快，数字系统的训练与验证也可能很慢。比特、原子与生命各自包含快过程和慢过程。

对确定的过程$i$，可以比较它的特征弛豫时间与外部驱动周期：
\begin{equation}
De_i=\frac{\tau_i}{T_{\text{驱动}}}.
\end{equation}
$\tau_i$与$T_{\text{驱动}}$采用相同时间单位，$De_i$无量纲。这一德博拉数来自材料研究；用于组织或生物干预时，要把“弛豫”落实为可以观察的恢复过程。

接收端恢复得比驱动快，相邻变化就可能分别被吸收；驱动快于恢复，尚未完成的响应就会重叠，历史开始进入当前状态。这个比值描述时间关系，稳定性还取决于增益、延迟、反馈方向、储备和执行能力。

仓库可以每隔几秒改变路线指令，人员训练与布局调整却要更久。电网可以迅速处理电气扰动，增建容量仍需多年。药物可以很快作用于靶点，长期效果与风险还在形成。项目真正需要识别的是这些时钟及其连接方式。

\begin{table}[ht]
\centering
\caption{快过程、慢接收端与可能积累的负担}
\label{tab:ch4-three-media}
{\small
\begin{tabularx}{\textwidth}{@{}l >{\RaggedRight\arraybackslash}X >{\RaggedRight\arraybackslash}X >{\RaggedRight\arraybackslash}X@{}}
\toprule
项目层 & 相对快的动作 & 较慢的接收过程 & 可能积累的负担\\
\midrule
数字 & 生成设计、提交改动 & 集成、审查与验证 & 未验证变更、版本债务\\
工业 & 修改订单与调度 & 建造设备、形成产能 & 积压订单、长期承诺\\
生命 & 启动干预 & 适应、恢复与长期观察 & 未解决暴露、风险与随访需求\\
\bottomrule
\end{tabularx}}
\end{table}

单一时钟往往不是好设计。局部快回路处理扰动，上层慢回路调整目标，更长周期检验容量与方向。分层、缓冲与明确接口，让不同速率能够共存。

\subsection{保住可以恢复的运行区间}

长期运行的系统需要一个区间：服务能够继续，扰动能够吸收，偏离之后仍有路径回来。流态的比喻由此获得用途：观察有序通过、波动放大与关键连接失效。实际判据要落在队列长度、恢复时间、共同失效、储备与关键服务连续性上。

局部改善可能把工作点推向边界。交期缩短，留给纠错的时间减少；库存下降，缓冲变薄；政策改动太密，证据尚未形成就进入下一轮。这些措施也能改进原本低效的过程。区别在于它们改变了怎样的完整回路。

第一条路径是配速。接收端确实需要时间，就拉长对应决策周期；测量与纠偏可以加快，就提高回路响应能力。快慢的价值取决于行动能否被接住。

第二条路径是分流。库存、储备、分阶段上线、模块边界与局部自治，都能隔开一部分速率差。缓冲容量应按扰动期间“流入超过流出”的累计量估计。可以补充的储备，与只把冲击推迟一次的等待，具有不同动力学。

第三条路径是维护稳定性所依赖的功能。验证产生证据，审计暴露问题，维护保存能力，责任安排连接行动与救济。它们的价值来自实际结果。能以更少步骤完成同一功能，就是进步；留下更长流程却没有留下功能，仍然只是成本。

把可靠性、缓冲和恢复纳入目标，优化也会选择这些投入。制度问题出现在支付与收益分离的地方：维护预算记在一个部门，避免的中断却惠及整个网络；设备寿命跨过数任管理者，今天的成绩可能在别人任内结账。

\key{必须有人跨过慢周期持有这笔账。}记录、预算和责任需要足够连续，才能把当前收益与后来的支出接在一起。生态影响可能超过许可周期，技能形成可能超过产品周期，公共设施可能超过任期。系统的时间跨度，应进入决策权的安排。

本章横跨三种介质，采用的是同一套核算纪律：身体怎样吸收接触误差，机器怎样减少或转移残差，车辆怎样改变能量流，运营网络怎样使用资产，飞行怎样缩短路径又建立空域秩序。成本可以净下降。完整边界让我们看见下降发生在哪里，剩余的时间、风险和费用由谁承担。

当符号行动变便宜，这套纪律面对另一种困难。方案、凭证、图像与身份，可以生成得比核实它们更快。对象、证据与承责者一旦错开，行动就可能取得效果，却找不到与效果相称的反馈。下一章从这里开始。

\furtherreading{%
Rolf Pfeifer and Josh Bongard，\textit{How the Body Shapes the Way We Think: A New View of Intelligence}，MIT Press，2006\\
Daniela Rus and Michael T. Tolley，“Design, Fabrication and Control of Soft Robots”，\textit{Nature}，2015\\
Richard C. Lewontin，\textit{The Triple Helix: Gene, Organism, and Environment}，Harvard University Press，2000\\
Evelyn Fox Keller，\textit{The Century of the Gene}，Harvard University Press，2000\\
Herbert A. Simon，\textit{The Sciences of the Artificial}，MIT Press，1969\\
Nancy G. Leveson，\textit{Engineering a Safer World: Systems Thinking Applied to Safety}，MIT Press，2011\\
Vaclav Smil，\textit{Energy and Civilization: A History}，MIT Press，2017\\
Jarrett Walker，\textit{Human Transit: How Clearer Thinking about Public Transit Can Enrich Our Communities and Our Lives}，Island Press，2012}


% <<<<<<<<<< parts/book-ch04.tex <<<<<<<<<<


% >>>>>>>>>> parts/book-ch05.tex >>>>>>>>>>
%% ===================== 本章参考文献 =====================
\printbibliography[heading=subbibliography,title=参考文献]
